\selectlanguage{english}
\chapter{System Design, Architecture, and Results}
\label{chap:design}
\renewcommand{\chaptername}{Chapter}

\section*{Introduction}
\addcontentsline{toc}{section}{\textbf{Introduction}}

This chapter presents the design, architecture, and preliminary results of the
proposed GEO measurement pipeline. It establishes the research approach and core
design decisions, then describes each component in detail: the dynamic model
registry, and the four specialized agents with their internal workflows. The chapter
concludes with preliminary evaluation results and a discussion of observed
limitations.

% ─────────────────────────────────────────────────────────────────────────────
\section{Research Approach and Design Philosophy}
% ─────────────────────────────────────────────────────────────────────────────

This project follows a \textbf{design science research} methodology
\cite{hevner2004design}: the primary contribution is a working software artifact
evaluated against the problem it was designed to solve. The approach involves
iterative construction and empirical evaluation --- each component is built, tested
under real API conditions, measured for output quality, and refined before the next
is designed. Three principles guided all architectural decisions.

\textbf{Autonomy over prescription:} the system makes its own decisions about which
model to use, which data source to query, and how to recover from failures ---
without human intervention at runtime. This ruled out fixed execution graphs in
favor of an orchestrator that reasons about the current state and selects the next
action dynamically.

\textbf{Resilience over speed:} the pipeline operates under constrained conditions
--- free-tier API quotas, variable rate limits, and tool servers that may time out.
The system must checkpoint progress and resume without data loss, motivating the
per-model state tracking in the model registry.

\textbf{Reproducibility over convenience:} all outputs are saved in structured CSV
files with fixed schemas. Intermediate outputs from each agent are preserved
separately so any stage can be re-run without repeating earlier stages.

% ─────────────────────────────────────────────────────────────────────────────
\section{System Overview}
% ─────────────────────────────────────────────────────────────────────────────

The pipeline is organized into four sequential stages, each handled by a dedicated
agent, coordinated by an orchestration layer, and supported by a dynamic model
registry shared across all agents. Figure~\ref{fig:system_overview} presents the
full system architecture.

\begin{figure}[H]
\centering
\begin{tikzpicture}[
  agent/.style={
    rectangle, rounded corners=5pt,
    draw=blue!65, fill=blue!8,
    text width=3.4cm, align=center,
    font=\small\bfseries, minimum height=1.0cm},
  orch/.style={
    rectangle, rounded corners=6pt,
    draw=violet!70, fill=violet!8,
    text width=2.4cm, align=center,
    font=\small\bfseries, minimum height=1.0cm},
  csv/.style={
    rectangle, rounded corners=3pt,
    draw=gray!55, fill=gray!5,
    text width=3.4cm, align=center,
    font=\scriptsize\itshape, minimum height=0.85cm},
  registry/.style={
    rectangle, rounded corners=5pt,
    draw=orange!70, fill=orange!8,
    text width=3.2cm, align=center,
    font=\small\bfseries, minimum height=0.9cm},
  provider/.style={
    rectangle, rounded corners=3pt,
    draw=red!40, fill=red!5,
    text width=2.8cm, align=center,
    font=\scriptsize, minimum height=0.6cm},
  mcp/.style={
    rectangle, rounded corners=3pt,
    draw=teal!60, fill=teal!7,
    text width=3.0cm, align=center,
    font=\scriptsize, minimum height=0.6cm},
  %% Arrow styles
  arrow/.style={-{Stealth[length=5pt]}, thick},
  darrow/.style={-{Stealth[length=4pt]}, thick, dashed, gray!60},
  oarrow/.style={-{Stealth[length=4pt]}, thick, orange!70},
  tarrow/.style={-{Stealth[length=4pt]}, thick, teal!65},
  earrow/.style={-{Stealth[length=4pt]}, thick, violet!65},
]

%% ═══════════════════════════════════════════════════════════════════════════
%% COLUMN POSITIONS
%%   Orchestrator  : x = -4.5   (left coordinator column)
%%   Agent column  : x =  0      (span: -1.7 to +1.7)
%%   CSV column    : x =  5.5    (span:  3.8 to  7.2)
%%   Routing lane  : x =  3.0    (gap between agent.east=1.7 and csv.west=3.8)
%%   Right column  : x =  9.8    (registry, providers, MCP)
%% ROW POSITIONS: y = 0, -3.5, -7.0, -10.5
%% ═══════════════════════════════════════════════════════════════════════════

%% ── Orchestrator (left column, vertically centered) ─────────────────────────
\node[orch] (orch) at (-4.5, -5.25) {Orchestrator\\{\footnotesize\normalfont orchestrator.py}};

%% ── Agents ──────────────────────────────────────────────────────────────────
\node[agent] (a0) at (0,    0)    {Agent 0\\{\footnotesize\normalfont Intent \& Prompts}};
\node[agent] (a1) at (0,   -3.5) {Agent 1\\{\footnotesize\normalfont Query \& Extract}};
\node[agent] (a2) at (0,   -7.0) {Agent 2\\{\footnotesize\normalfont Web Research}};
\node[agent] (a3) at (0,  -10.5) {Agent 3\\{\footnotesize\normalfont Merge \& Score}};

%% ── CSV checkpoints ─────────────────────────────────────────────────────────
\node[csv] (c0) at (5.5,    0)   {\texttt{prompt\_set.csv}\\{\tiny intents, prompts, language}};
\node[csv] (c1) at (5.5,  -3.5) {\texttt{entity\_features\_global.csv}\\{\tiny entities, mention\_rate, stability}};
\node[csv] (c2) at (5.5,  -7.0) {\texttt{web\_features.csv}\\{\tiny ratings, social, wikipedia}};
\node[csv] (c3) at (5.5, -10.5) {\texttt{enriched\_entities.csv}\\{\tiny all features + geo\_score}};

%% ── Model Registry + Providers (right column, rows 0–1 area) ────────────────
\node[registry] (reg)  at (9.8, -1.0)  {Dynamic Model\\Registry};
\node[provider] (groq) at (9.8, -2.5)  {Groq API};
\node[provider] (or)   at (9.8, -3.6)  {OpenRouter};
\node[provider] (mist) at (9.8, -4.7)  {Mistral API};

%% ── MCP Tool Servers (right column, row 2 area) ─────────────────────────────
\node[mcp] (mcp1) at (9.8, -6.2)  {Search Server\\{\tiny Maps · TA · DDG}};
\node[mcp] (mcp2) at (9.8, -7.2)  {Scrape Server\\{\tiny Instagram · Apify}};
\node[mcp] (mcp3) at (9.8, -8.2)  {Wiki Server\\{\tiny Wikipedia EN/FR}};

%% ═══════════════════════════════════════════════════════════════════════════
%% ARROWS
%% ═══════════════════════════════════════════════════════════════════════════

%% 0. Orchestrator → Agents  (sequential calls — fixed pipeline order)
%%    Route: orch.east → horizontal to x=-2.1, then bend down/up to each agent.west
\draw[earrow] (orch.east) -| (a0.west)
  node[near end, left, font=\scriptsize, violet!70] {(1)};
\draw[earrow] (orch.east) -| (a1.west)
  node[near end, left, font=\scriptsize, violet!70] {(2)};
\draw[earrow] (orch.east) -| (a2.west)
  node[near end, left, font=\scriptsize, violet!70] {(3)};
\draw[earrow] (orch.east) -| (a3.west)
  node[near end, left, font=\scriptsize, violet!70] {(4)};

%% 1. Agent → CSV  (horizontal writes, same row — always clear)
\draw[arrow] (a0.east) -- (c0.west)
  node[midway, above, font=\scriptsize] {writes};
\draw[arrow] (a1.east) -- (c1.west)
  node[midway, above, font=\scriptsize] {writes};
\draw[arrow] (a2.east) -- (c2.west)
  node[midway, above, font=\scriptsize] {writes};
\draw[arrow] (a3.east) -- (c3.west)
  node[midway, above, font=\scriptsize] {writes};

%% 2. CSV → next Agent  (sequential data feeds)
%%    Route: csv.south ↓ into the inter-row gap → left → agent.east
%%    All horizontals sit in clear inter-row gaps (no nodes there)

%% c0 → a1  (gap at y ≈ -1.75)
\draw[darrow]
  (c0.south) -- ++(0,-1.3) -| (a1.east)
  node[near start, right, font=\scriptsize, gray!70] {reads};

%% c1 → a2  (gap at y ≈ -5.25)
\draw[darrow]
  (c1.south) -- ++(0,-1.3) -| (a2.east);

%% c2 → a3  (gap at y ≈ -8.75)
\draw[darrow]
  (c2.south) -- ++(0,-1.3) -| (a3.east);

%% 3. c1 → a3  (Agent 3 also reads Agent 1 output — bypass via routing lane)
%%    Routing lane at x=3.0: between agent.east (1.7) and csv.west (3.8) — no nodes
\draw[darrow, gray!45]
  (c1.west)
    -- (3.0, -3.5)   %% enter the routing lane at row-1 level
    -- (3.0, -10.5)  %% descend to row-3 level (no nodes at x=3.0)
    -- (a3.east)
  node[near start, left, font=\scriptsize, gray!55] {also reads};

%% 4. Registry → LLM Providers  (downward within right column — always clear)
\draw[oarrow] (reg.south)  -- (groq.north);
\draw[oarrow] (groq.south) -- (or.north);
\draw[oarrow] (or.south)   -- (mist.north);

%% 5. Registry → Agent 1 and Agent 2  (model selection)
%%    Route: reg.west → routing lane (x=3.0) → agent.east
%%    Horizontal leg stays in the inter-row gap; vertical leg in clear lane

%% reg → a1: horizontal at y=-1.0 (between row-0 and row-1, no csv nodes there)
\draw[oarrow]
  (reg.west)
    -- (3.1, -1.0)   %% reach the routing lane at a clear y level
    -- (3.1, -3.5)   %% descend to row-1 in the lane
    -- (a1.east)
  node[near end, above, font=\scriptsize, orange!80] {model};

%% reg → a2: horizontal at y=-1.0 same entry, continues down to row-2
\draw[oarrow]
  (reg.west)
    -- (2.9, -1.0)   %% slightly offset to avoid exact overlap with reg→a1
    -- (2.9, -7.0)   %% descend to row-2 in the lane
    -- (a2.east);

%% 6. Agent 2 ↔ MCP Tool Servers
%%    Route: a2.east → ++(1.5,0) — reaching x=1.7+1.5=3.2 (in the lane)
%%    then right to MCP, staying above csv.south and below next row
\draw[tarrow]
  (a2.east) -- ++(2.0,0) |- (mcp1.west)
  node[near start, above, font=\scriptsize, teal!70] {tool calls};
\draw[tarrow] (a2.east) -- ++(2.0,0) |- (mcp2.west);
\draw[tarrow] (a2.east) -- ++(2.0,0) |- (mcp3.west);

%% ── Grouping boxes ──────────────────────────────────────────────────────────
\node[draw=teal!40, dashed, rounded corners=5pt,
      fit=(mcp1)(mcp2)(mcp3), inner sep=4pt,
      label={[font=\scriptsize\itshape, teal!60]above:MCP Tool Servers}] {};

\node[draw=red!30, dashed, rounded corners=5pt,
      fit=(groq)(or)(mist), inner sep=4pt,
      label={[font=\scriptsize\itshape, red!40]above:LLM Providers}] {};

\end{tikzpicture}
\caption{GEO pipeline data flow. The orchestrator (\texttt{orchestrator.py})
drives the four agents in fixed sequential order (1)–(4). Each agent reads its
input from the upstream CSV checkpoint and writes its output to the next.
Agent~3 merges both Agent~1 and Agent~2 outputs using fuzzy entity matching,
then applies LLM-based deduplication and entity triage.
Agents~1 and~2 query the dynamic model registry for the best available
(model, provider, key) triple at each inference call. Agent~2 performs all
web data collection exclusively through the three MCP tool servers.}
\label{fig:system_overview}
\end{figure}

% ─────────────────────────────────────────────────────────────────────────────
\section{Dynamic Model Registry}
% ─────────────────────────────────────────────────────────────────────────────

Operating across multiple free-tier LLM providers requires managing heterogeneous
and frequently binding quota constraints. The \textbf{dynamic model registry}
maintains a continuously updated state for every available model slot and selects
the best available option autonomously at each inference call --- eliminating the
need for hardcoded fallback chains.

At pipeline startup, the registry queries the model listing APIs of Groq, OpenRouter,
and Mistral, constructing a pool of slots identified by
\texttt{provider:keyIndex/model\_id}. Each slot is assigned a capability tier
(Table~\ref{tab:model_tiers}) and maintains three runtime states: \textbf{TPD
exhaustion} (daily quota hit --- permanent, persisted to disk across restarts),
\textbf{TPM block} (per-minute rate limit --- temporary, cleared when the retry
window expires), and \textbf{learned token capacity} (set on 413 errors --- used to
pre-filter the model for oversized prompts, eliminating repeated failures without
manual configuration).

\begin{table}[H]
\centering
\caption{Model capability tiers used for selection ranking.}
\label{tab:model_tiers}
\begin{tabular}{|c|p{3.2cm}|p{8.3cm}|}
\hline
\textbf{Tier} & \textbf{Quality Level} & \textbf{Identifying Keywords} \\
\hline
0 & Highest   & 70b, 72b, 120b, llama-4, qwen3-32b, claude-3.5, deepseek-r1 \\
\hline
1 & Medium    & 32b, 27b, 14b, gemini-flash, mistral-small, claude-3-haiku \\
\hline
2 & Lightweight & 8b, 7b, 3b, instant, mini, nano \\
\hline
\end{tabular}
\end{table}

The \texttt{select(preferred, estimated\_tokens)} method ranks candidates by tier,
provider, and key index --- but the pool it operates on evolves autonomously at
runtime. Three state transitions shape it without human intervention:
\texttt{mark\_tpd\_exhausted()} permanently removes a slot when its daily quota is
hit (persisted across restarts), \texttt{mark\_tpm()} blocks a slot until its
rate-limit window expires, and \texttt{mark\_too\_large()} records a learned
token-capacity ceiling on 413 errors so oversized prompts skip that model
automatically. The pipeline thus always selects the best \textit{currently usable}
model for the prompt at hand --- \textbf{adaptive autonomous selection} driven
entirely by the system's own API observations. The error-classification logic that
triggers these transitions lives in the Supervisor Agent
(Section~\ref{sec:supervisor}), described next.

% ─────────────────────────────────────────────────────────────────────────────
\section{Pipeline Agents}
% ─────────────────────────────────────────────────────────────────────────────

\subsection{Agent 0: Intent Discovery and Prompt Generation}

A single query formulation such as ``best restaurants in Tunis'' would produce
results sensitive to the experimenter's specific phrasing. Agent~0 addresses this
by autonomously generating a diverse prompt set covering multiple user intents and
linguistic variants, ensuring that downstream visibility metrics reflect robust
entity representation across a broad query space.

The generation pipeline consists of three steps. \textbf{Step 1 --- Intent
discovery:} given a domain description and target language, an LLM identifies
the distinct query intents a real user might have (e.g.\ recommendation requests,
occasion-based searches, comparison queries, information lookup). The number of
intents is controlled by the \texttt{n\_intents} parameter; the result is a
validated JSON taxonomy stored in the pipeline state. \textbf{Step 2 --- Prompt
generation:} for each intent and each target language (French and Arabic in this study),
\texttt{n\_variants} distinct formulations are generated using varied vocabulary,
sentence structure, and specificity level, simulating natural multilingual user
variation. \textbf{Step 3 --- Self-reflection:} the LLM evaluates the \textit{entire} prompt
set against quality criteria: intent diversity, linguistic naturalness, and a
GEO-specific constraint requiring that at least 60\% of prompts force the AI to
name specific establishments (using superlatives, rankings, or explicit citation
requests). If the set fails (\texttt{proceed: false}), the whole generation is
restarted with the missing intents added to \texttt{n\_intents}; the loop runs at
most \texttt{max\_reflection\_loops} times (default~2). The final set is saved to
\texttt{prompt\_set.csv}.

\begin{figure}[H]
\centering
\begin{tikzpicture}[
  box/.style={rectangle, rounded corners=4pt, draw=blue!60, fill=blue!8,
              text width=3.8cm, align=center, font=\small, minimum height=1.0cm},
  decision/.style={rectangle, rounded corners=4pt, draw=orange!70, fill=orange!8,
                   text width=3.0cm, align=center, font=\scriptsize,
                   minimum height=0.8cm},
  arrow/.style={-{Stealth[length=5pt]}, thick},
  io/.style={rectangle, rounded corners=4pt, draw=green!60!black, fill=green!7,
             text width=3.4cm, align=center, font=\small, minimum height=0.8cm},
]
\node[io]       (input)   at (0, 0)    {Domain + Language\\+ Parameters};
\node[box]      (intents) at (0, -2.0) {Step 1 --- Intent Discovery\\{\footnotesize LLM intent taxonomy}};
\node[box]      (prompts) at (0, -4.0) {Step 2 --- Prompt Generation\\{\footnotesize n\_variants per intent}};
\node[box]      (reflect) at (0, -6.0) {Step 3 --- Self-Reflection\\{\footnotesize quality scoring}};
\node[decision] (check)   at (0, -8.0) {All prompts pass\\quality threshold?};
\node[io]       (output)  at (5.5, -8.0) {\texttt{prompt\_set.csv}};
\draw[arrow] (input)   -- (intents);
\draw[arrow] (intents) -- (prompts);
\draw[arrow] (prompts) -- (reflect);
\draw[arrow] (reflect) -- (check);
\draw[arrow] (check.east) -- (output.west)
  node[midway, above, font=\scriptsize] {Yes};
\draw[arrow] (check.west) -- ++(-2.5, 0) |- (prompts.west)
  node[near start, left, font=\scriptsize, align=center] {No\\regenerate};
\end{tikzpicture}
\caption{Agent~0 workflow: intent discovery, prompt generation, and self-reflection
quality control. Low-quality prompts trigger a targeted regeneration loop.}
\label{fig:agent0_flow}
\end{figure}

\subsection{Agent 1: LLM Querying and Entity Extraction}

Agent~1 is responsible for the core GEO measurement task: submitting the generated
prompts to multiple LLMs across independent runs, extracting the named entities
mentioned in each response, and computing visibility metrics that quantify how
consistently each entity appears across the full prompt, model, and run space.

\textbf{Step 2 --- Multi-model querying.} Each prompt is submitted to every
model slot in the \texttt{QUERY\_MODELS} configuration, repeated \texttt{n\_runs}
times per slot. This cross-model, multi-run design is the core of GEO measurement:
entity visibility is only meaningful when aggregated across diverse LLMs and
phrasing variants, not a single model call. All responses are appended
progressively to \texttt{raw\_responses.csv} with resume support.

\textbf{Step 3 --- Zero-shot entity extraction.} A dedicated Mistral extractor
(fixed model, temperature~0) processes each raw response with a detailed system
prompt specifying domain scope (Tunisian establishments only), normalization rules
(lowercase, strip generic prefixes), and output format (flat JSON string array).
The extractor returns entity names only --- no confidence scores --- along with the
text block describing each entity for downstream context.

\textbf{Step 4 --- Enrichment.} For each extracted entity, two structural features
are computed: \textbf{ranking position} (the order in which the entity first
appeared in the response, extracted via a second LLM call over the ranked list) and
\textbf{description length in tokens} (estimated from the character ratio of the
entity's description block to the full response).

\textbf{Step 5 --- Five-phase cleaning pipeline.}
(A)~\textbf{Hard pre-filtering} drops empty, too-short ($<$3 chars), numeric, and
singleton entries (mention count $<$ 2) --- singletons are noise, not signal.
(B)~\textbf{Fuzzy clustering} (RapidFuzz, threshold~88) groups high-similarity
names; the most frequent variant in each cluster becomes the canonical form.
(C)~\textbf{LLM arbitration} resolves ambiguous pairs (similarity 75--87) and
simultaneously validates every entity: entities that are generic descriptors, wrong
geography, or lack Tunisian context are invalidated.
(D)~\textbf{Self-reflection} reviews all cleaning decisions and immediately applies
corrections --- undoing incorrect merges or invalidations flagged as suspicious.
(E)~\textbf{Mapping application} writes the final canonical names, quality flags,
and null-rank penalties to \texttt{clean\_entities.csv}.

\textbf{Step 6 --- GEO metric computation.} Three visibility metrics are computed
per canonical entity across all responses: \textbf{mention rate} (fraction of
valid responses in which the entity appeared), \textbf{stability score}
($\text{mention\_rate} \times \frac{1}{1+\text{rank\_variance}}$, combining
frequency with positional consistency), and a \textbf{consistency label}
(\textit{stable} $\geq$~0.75, \textit{variable} $\geq$~0.40, \textit{unstable}
otherwise). Entities with fewer than 2 mentions are excluded at Phase~A; the
stability threshold for Agent~2 eligibility is applied by the orchestrator.

\begin{figure}[H]
\centering
\begin{tikzpicture}[
  box/.style={rectangle, rounded corners=4pt, draw=blue!60, fill=blue!8,
              text width=4.0cm, align=center, font=\small, minimum height=1.0cm},
  arrow/.style={-{Stealth[length=5pt]}, thick},
  io/.style={rectangle, rounded corners=4pt, draw=green!60!black, fill=green!7,
             text width=3.4cm, align=center, font=\small, minimum height=0.8cm},
]
\node[io]  (input)   at (0, 0)     {\texttt{prompt\_set.csv}};
\node[box] (query)   at (0, -2.0)  {Step 2 --- Multi-model Querying\\{\footnotesize QUERY\_MODELS $\times$ n\_runs}};
\node[box] (extract) at (0, -4.0)  {Step 3 --- Entity Extraction\\{\footnotesize Mistral, flat JSON, names only}};
\node[box] (enrich)  at (0, -6.0)  {Step 4 --- Enrichment\\{\footnotesize ranking position + desc.\ length}};
\node[box] (clean)   at (0, -8.0)  {Step 5 --- Cleaning (A--E)\\{\footnotesize hard rules · fuzzy · LLM · reflect}};
\node[box] (metrics) at (0, -10.0) {Step 6 --- GEO Metrics\\{\footnotesize mention rate, stability score}};
\node[io]  (output)  at (0, -12.0) {\texttt{entity\_features\_global.csv}};
\draw[arrow] (input)   -- (query);
\draw[arrow] (query)   -- (extract)
  node[midway, right, font=\scriptsize] {\texttt{raw\_responses.csv}};
\draw[arrow] (extract) -- (enrich)
  node[midway, right, font=\scriptsize] {\texttt{extracted\_entities.csv}};
\draw[arrow] (enrich)  -- (clean)
  node[midway, right, font=\scriptsize] {\texttt{enriched\_entities.csv}};
\draw[arrow] (clean)   -- (metrics)
  node[midway, right, font=\scriptsize] {\texttt{clean\_entities.csv}};
\draw[arrow] (metrics) -- (output);
\end{tikzpicture}
\caption{Agent~1 workflow: multi-model querying, entity extraction, enrichment,
five-phase cleaning, and GEO metric computation. Each intermediate output is
checkpointed to CSV for resume support.}
\label{fig:agent1_flow}
\end{figure}

\subsection{Agent 2: Web Presence Research}

Agent~2 is the most technically complex component. For each entity meeting the
visibility threshold, it operates as a fully autonomous ReAct agent over three
MCP tool servers --- at each step it produces a \textit{thought} about what
information has been gathered so far, selects a \textit{tool invocation} as its
\textit{action}, and incorporates the tool's response as an \textit{observation}
before deciding the next step.

The agent follows a \textbf{prescribed 8-step research protocol} defined in its
system prompt: (1)~\texttt{google\_maps\_search}, (2)~\texttt{wikipedia\_lookup},
(3)~\texttt{tripadvisor\_search}, (4)~\texttt{wikidata\_lookup}, then conditionally:
(5)~\texttt{extract\_website\_socials} if a website was found, (6)~\texttt{scrape\_instagram}
if an Instagram handle was found, (7)~\texttt{scrape\_facebook} if a Facebook handle
was found, and (8)~\texttt{ddg\_search} if no social presence was located. Steps 1--4
always execute in order; steps 5--8 depend on what earlier steps returned.
The ReAct loop provides adaptive handling of tool failures and observations, but the
sequence of data targets is fixed.

The four MCP servers expose the tools described in Table~\ref{tab:mcp_tools}. The
protocol collects data across six categories (Google Maps, TripAdvisor, Wikipedia,
Wikidata, Instagram, Facebook) and enforces geographic anchoring: every search query
must include ``Tunisia'' or a Tunisian city name, and any result referencing a
non-Tunisian geographic context is discarded and the search retried with more
specific terms.

\begin{table}[H]
\centering
\caption{MCP tool servers and their exposed tools.}
\label{tab:mcp_tools}
\begin{tabular}{|p{2.5cm}|p{3.5cm}|p{7.2cm}|}
\hline
\textbf{Server} & \textbf{Tool} & \textbf{Description} \\
\hline
\multirow{3}{*}{search\_server}
  & \texttt{google\_maps\_search}      & Rating, review count, address, phone, website. Rotates through a pool of SerpAPI keys. \\
\cline{2-3}
  & \texttt{tripadvisor\_search}       & Two-step: web search locates the TA page URL, then scraper extracts rating and ranking. \\
\cline{2-3}
  & \texttt{ddg\_search}               & General DuckDuckGo search returning titles, URLs, snippets. \\
\hline
\multirow{3}{*}{scrape\_server}
  & \texttt{scrape\_instagram}         & Follower count, post count, biography via Apify. Rotates through 5 Apify tokens. \\
\cline{2-3}
  & \texttt{scrape\_facebook}          & Page likes and average post engagement via Apify. \\
\cline{2-3}
  & \texttt{extract\_website\_socials} & Extracts Instagram/Facebook handles from a business website's HTML. \\
\hline
wiki\_server
  & \texttt{wikipedia\_lookup}         & Queries the Wikipedia API in English and French; returns boolean presence and URL. \\
\hline
enrichment\_server
  & \texttt{wikidata\_lookup}          & Queries the Wikidata API; returns entity type, country, founded date, and official website. \\
\hline
\end{tabular}
\end{table}

\subsubsection*{Supervisor Agent}
\label{sec:supervisor}

API errors are handled by the \textbf{Supervisor Agent} --- a rule-based error
classifier with an optional LLM escalation path --- that inspects each error message
and returns a structured recovery decision
with four possible actions: \texttt{switch\_model} (billing or quota exhaustion ---
mark slot TPD-exhausted), \texttt{wait\_retry} (rate limit --- wait the extracted
retry-after duration), \texttt{retry} (transient 5xx --- retry immediately), or
\texttt{skip} (unrecoverable after 4 attempts). Internally, the supervisor first
attempts to call a lightweight LLM on the error text; if that call fails or times
out, it falls back to a deterministic \texttt{if/elif} chain that encodes the same
classification rules explicitly. In practice the LLM and the fallback produce
identical decisions for the structured error codes that providers return --- the LLM
path adds resilience for novel or ambiguous error messages, not a fundamentally
different decision process. Each successfully processed entity is immediately
checkpointed to \texttt{web\_features.csv}.

\begin{figure}[H]
\centering
\begin{tikzpicture}[
  box/.style={rectangle, rounded corners=4pt, draw=blue!60, fill=blue!8,
              text width=3.6cm, align=center, font=\small, minimum height=1.0cm},
  supervisor/.style={rectangle, rounded corners=4pt, draw=red!50, fill=red!6,
              text width=3.2cm, align=center, font=\small, minimum height=1.0cm},
  mcp/.style={rectangle, rounded corners=4pt, draw=teal!60, fill=teal!7,
              text width=3.2cm, align=center, font=\small, minimum height=1.0cm},
  io/.style={rectangle, rounded corners=4pt, draw=green!60!black, fill=green!7,
             text width=3.2cm, align=center, font=\small, minimum height=0.8cm},
  arrow/.style={-{Stealth[length=5pt]}, thick},
]
\node[io]         (input)  at (0, 0)     {Entity Name + Context};
\node[box]        (think)  at (0, -2.0)  {Thought\\Reason about gathered data};
\node[box]        (act)    at (0, -4.0)  {Action\\Select \& invoke MCP tool};
\node[mcp]        (mcp)    at (5.0, -4.0){MCP Tool Servers\\Search · Scrape · Wiki · Enrichment};
\node[box]        (obs)    at (0, -6.0)  {Observation\\Parse \& validate result};
\node[box]        (check)  at (0, -8.0)  {Sufficient data\\collected?};
\node[supervisor] (sup)    at (5.0, -6.5){Supervisor Agent\\Classify error\\Return action};
\node[io]         (output) at (0, -10.5) {\texttt{web\_features.csv}};
\draw[arrow] (input)  -- (think);
\draw[arrow] (think)  -- (act);
\draw[arrow] (act.east)  -- (mcp.west)  node[midway, above, font=\scriptsize] {tool call};
\draw[arrow] (mcp.west)  -- (obs.east)  node[midway, below, font=\scriptsize] {result};
\draw[arrow] (obs)    -- (check);
\draw[arrow] (check.west) -- ++(-2.2, 0) |- (think.west)
  node[near start, left, font=\scriptsize, align=center] {No\\next step};
\draw[arrow] (check.south) -- (output)  node[midway, right, font=\scriptsize] {Yes};
\draw[arrow, red!50, dashed] (obs.east)  -- (sup.west)
  node[midway, above, font=\scriptsize] {on error};
\draw[arrow, red!50, dashed] (sup.north) -- ++(0, 0.6) -| (act.east)
  node[near start, right, font=\scriptsize, align=center] {retry /\\switch model};
\end{tikzpicture}
\caption{Agent~2 ReAct processing loop. The agent iterates
Thought--Action--Observation cycles over MCP tool servers until sufficient web
presence data has been collected. All errors are handled by the Supervisor Agent,
which reasons about the failure type and returns a structured recovery action
without terminating the loop.}
\label{fig:agent2_flow}
\end{figure}

\subsection{Agent 3: Merge, Deduplication, and Agentic Triage}

Agent~3 is the final stage of the pipeline. It joins the LLM visibility metrics
from Agent~1 with the web presence features from Agent~2 using a three-stage fuzzy
matching strategy (exact normalized match, substring containment, then word overlap),
with Agent~1 as the source of truth so that every extracted entity is retained.

Two LLM-based reasoning steps then clean the merged dataset. First,
\textbf{deduplication}: entities are grouped by shared content words (stop words
such as \textit{le, la, de, el, dar} removed) and each candidate group is submitted
to the LLM to decide whether the rows refer to the same place. Confirmed duplicates
are merged --- mention counts summed, the richer web-data row kept.
Second, \textbf{triage}: entities for which Agent~2 found no web evidence are
classified by the LLM as \textit{retry} (real restaurant, missed by Agent~2),
\textit{hallucination}, \textit{non\_restaurant}, or \textit{generic}.
Non-viable entities are removed and saved to \texttt{rejected\_entities.csv};
\textit{retry} entities are written to the merge report so the orchestrator can
re-invoke Agent~2 on them.

Finally, Agent~3 computes a \textbf{data-completeness score} (fraction of ten key
fields that are non-null) and assigns a \textbf{unified quality} label
(\textit{complete}, \textit{partial}, \textit{sparse}, or \textit{agent2\_failed})
to each row. The result is saved to \texttt{unified\_features.csv}.
Table~\ref{tab:output_schema} lists the key output fields.

\begin{table}[H]
\centering
\caption{Key fields in the final output dataset (\texttt{unified\_features.csv}).}
\label{tab:output_schema}
\begin{tabular}{|p{3.8cm}|p{2.2cm}|p{7.2cm}|}
\hline
\textbf{Field} & \textbf{Source} & \textbf{Description} \\
\hline
\texttt{canonical\_entity}   & Agent 1 & Normalized entity name \\
\texttt{mention\_rate}       & Agent 1 & Fraction of valid responses with a mention \\
\texttt{stability\_score}    & Agent 1 & $\text{mention\_rate} \times \frac{1}{1+\text{rank\_variance}}$ \\
\texttt{consistency\_label}  & Agent 1 & stable / variable / unstable \\
\texttt{gm\_rating}          & Agent 2 & Google Maps rating (0--5) \\
\texttt{gm\_review\_count}   & Agent 2 & Google Maps number of reviews \\
\texttt{ta\_rating}          & Agent 2 & TripAdvisor rating \\
\texttt{ta\_ranking}         & Agent 2 & TripAdvisor ranking position \\
\texttt{has\_wikipedia}      & Agent 2 & Boolean Wikipedia presence \\
\texttt{ig\_followers}       & Agent 2 & Instagram follower count \\
\texttt{overall\_confidence} & Agent 2 & Agent-assigned data quality score (0--1) \\
\texttt{data\_completeness}  & Agent 3 & Fraction of key fields non-null (0--1) \\
\texttt{unified\_quality}    & Agent 3 & complete / partial / sparse / agent2\_failed \\
\texttt{triage\_decision}    & Agent 3 & retry / hallucination / non\_restaurant / generic (failed entities only) \\
\texttt{triage\_reason}      & Agent 3 & LLM justification for the triage decision \\
\hline
\end{tabular}
\end{table}

% ─────────────────────────────────────────────────────────────────────────────
\section{Results}
% ─────────────────────────────────────────────────────────────────────────────

The pipeline was applied to the Tunisian restaurant domain. Agent~0 generated
\textbf{30 prompts} covering \textbf{6 distinct user intents} (top recommendation,
tourist reviews, comparative, authentic cuisine, special occasion, location-based)
with 5~variants each in French. Agent~1 submitted each prompt to
\textbf{3 model slots}, repeated 5~times per slot, yielding \textbf{450 total
responses}. From these, 2{,}020 raw entity extraction records (831 unique surface
forms) were produced, which the five-phase cleaning pipeline consolidated into
\textbf{234 canonical entities}.

Of the 234 entities, \textbf{213 met the eligibility threshold} (minimum mention
rate) and were submitted to Agent~2 for web enrichment. Agent~3 then applied
LLM-based deduplication to collapse variant spellings of the same establishment
(e.g., \textit{Dar El Jeld} vs.\ \textit{Dar el Jeld Restaurant}) and
LLM-based triage to remove non-viable entries: 4~entities were rejected
(1~hallucination, 1~non-restaurant venue, 2~over-generic names). The final
dataset contains \textbf{182 enriched entities}. Table~\ref{tab:pipeline_results}
summarizes the key metrics from the actual output files.

\begin{table}[H]
\centering
\caption{Pipeline results on the Tunisian restaurant domain.}
\label{tab:pipeline_results}
\begin{tabular}{|p{7cm}|r|}
\hline
\textbf{Metric} & \textbf{Value} \\
\hline
Prompts generated (Agent 0)                          & 30 \\
LLM responses collected (Agent 1)                    & 450 \\
Raw entity extraction records                        & 2{,}020 \\
Unique surface forms before cleaning                 & 831 \\
Canonical entities after cleaning                    & 234 \\
Entities submitted to Agent 2                        & 213 \\
Final entities after dedup + triage (Agent 3)        & 182 \\
\quad — Complete ($\geq$3 sources)                   & 75 \\
\quad — Partial (1--2 sources)                       & 84 \\
\quad — Failed (no sources recovered)                & 22 \\
Entities rejected by LLM triage                      & 4 \\
Google Maps coverage                                 & 85.2\% \\
TripAdvisor coverage                                 & 38.5\% \\
Wikipedia presence                                   & 0\% \\
Instagram presence                                   & 31.9\% \\
Facebook presence                                    & 28.0\% \\
Mean data-completeness score                         & 0.584 \\
Mean overall confidence score                        & 0.541 \\
\hline
\end{tabular}
\end{table}

The distribution of LLM visibility is strongly concentrated: the top 15\% of
entities (27 out of 182) accumulate 62\% of all 1{,}351 mention instances, while
the remaining 85\% share the other 38\%. The single most-visible establishment,
\textit{Dar El Jeld}, received 159 mentions across all prompts and model slots ---
more than twice the count of the second-ranked entity. All entities carry an
\textit{unstable} consistency label, which is expected given the limited prompt set
(30 prompts) and the high variance of mention rates at this scale. The zero
Wikipedia coverage confirms that none of the 182 entities has a dedicated Wikipedia
article, consistent with the local and informal character of the Tunisian restaurant
market. The moderate Instagram and Facebook presence (32\% and 28\% respectively)
reflects the partial digitalization of small hospitality businesses in Tunisia.

% ─────────────────────────────────────────────────────────────────────────────
\section{Limitations and Discussion}
% ─────────────────────────────────────────────────────────────────────────────

Despite the agentic design improvements, several limitations were encountered.

\textbf{Partial web enrichment.} Of the 182 final entities, 22 (12\%) remain in a
\textit{failed} state after one enrichment pass and a second targeted retry pass.
These are restaurants whose names are either too common (shared with many
unrelated venues), spelled inconsistently across sources, or genuinely absent from
indexed web pages. Improving recall for this tail would require more specific
search strategies, such as incorporating city-level geocoding into the SerpAPI
query or combining multiple web sources per entity.

\textbf{Geographic disambiguation.} A number of entity names extracted from LLM
responses refer to well-known establishments outside Tunisia --- for example,
\textit{Le Grand Véfour} (Paris) or \textit{La Mamounia} (Marrakech hotel) ---
that are far better represented in LLM training data than any Tunisian homonym.
Agent~3's LLM triage successfully identified and rejected 4 such non-viable
entities, but the root cause lies upstream: Agent~0 prompts should more
aggressively enforce geographic specificity to prevent the LLM from surfacing
non-Tunisian establishments in the first place.

\textbf{Social media data sparsity.} Only 31.9\% of the 182 entities carry an
identified Instagram handle and 28.0\% a Facebook page, reflecting the relatively
low digital formalization of small hospitality businesses in the Tunisian market.
This sparsity limits the social dimension of the GEO score for the majority of
entities, and may understate the actual social media presence of establishments
that operate under informal or unverified page names.

\textbf{Triage robustness and JSON fragility.} The LLM triage step in Agent~3
serializes entity names into JSON prompts. Entity names containing apostrophes
(e.g., \textit{la table d'or}) caused JSON parse failures in initial runs,
requiring a robust fallback parser (\texttt{\_parse\_json\_objects}) that extracts
individual JSON objects when the full array parse fails. This fragility is an
inherent risk of relying on LLM-generated structured output and should be
addressed by escaping entity names before injection or switching to a
tool-calling interface.

\textbf{Retry convergence.} The two-pass retry design (Agent~2 $\rightarrow$
Agent~3 triage $\rightarrow$ Agent~2 retry) leaves 21 entities in a pending retry
queue at the end of the observed execution. A fully converged pipeline would
require either a third pass or a dedicated loop termination criterion (e.g., no
improvement after two consecutive retries).

\textbf{API quota and rate-limit constraints.}
The pipeline relies on three categories of external APIs, each with distinct
failure modes. \textit{LLM inference APIs} (Groq, OpenRouter, Together AI, Mistral)
impose per-minute and per-day token quotas on free tiers; hitting these limits
mid-pipeline caused request timeouts and forced the dynamic model registry to
cycle through fallback models, sometimes exhausting all available slots before a
batch completed. \textit{SerpAPI}, used by Agent~2 for Google Maps and TripAdvisor
lookups, provides only 100 free queries per key; with 213 entities requiring
multiple search queries each, the five available keys were consumed in a single
execution pass, leaving a subset of entities without web data until the keys
reset. \textit{Wikidata and Instagram} imposed no hard quotas but required
progressive throttling (sleep intervals between requests) to avoid HTTP~429
responses. Collectively, these constraints forced pipeline execution to be split
across multiple sessions on different days, introducing a temporal inconsistency in
the collected data. A production deployment would require paid API tiers or
self-hosted inference to run the full pipeline in a single uninterrupted pass.

\subsection*{Toward a GEO-Based Recommendation System}

Beyond measurement, the ultimate objective of this project is to build a
\textbf{recommendation system for brand visibility in generative AI}.
The pipeline described in this chapter constitutes the data acquisition and
diagnostic layer: it quantifies, for each restaurant, how visible it currently is
to LLM-based search engines and how complete its digital footprint is across the
web. The next step is to exploit this dataset to generate \textit{actionable
recommendations} --- identifying, for each brand, the specific levers that would
increase its likelihood of being cited by an AI assistant.

Concretely, a low \texttt{mention\_rate} combined with high \texttt{gm\_review\_count}
suggests that the brand has real-world credibility but lacks the textual signal
(structured web mentions, linked data, editorial coverage) that LLMs rely on
during pre-training. Conversely, a brand with high LLM visibility but low
TripAdvisor coverage may benefit from consolidating its review presence to
reinforce consistency across sources. The GEO score synthesizes these dimensions
into a single prioritization signal, enabling the recommendation engine to rank
optimization actions by expected impact on generative visibility.

This framing --- measuring current AI visibility, diagnosing its root causes,
and prescribing targeted improvements --- positions the system as a practical
tool for restaurant owners and marketers who wish to compete in the emerging
landscape where AI-generated responses, rather than ranked search results, drive
consumer discovery.

% ─────────────────────────────────────────────────────────────────────────────
\section*{Conclusion}
\addcontentsline{toc}{section}{\textbf{Conclusion}}
% ─────────────────────────────────────────────────────────────────────────────

This chapter has presented the full design, architecture, and results of the GEO
measurement pipeline. The four agents --- each with a well-defined input, processing
workflow, and structured output --- collectively realize the three design principles:
autonomy through supervisor-based error recovery, registry-driven model selection,
and LLM-based deduplication and triage; resilience through progressive checkpointing,
fuzzy entity matching, and a two-pass retry loop; and reproducibility through fixed
CSV schemas and decoupled outputs. Applied to the Tunisian restaurant domain, the
pipeline produced a final enriched dataset of \textbf{182 entities} with 85.2\%
Google Maps coverage and a mean data-completeness of 0.584, confirming both a
strong concentration of LLM visibility (top 15\% of entities attract 62\% of
mentions) and the practical challenges of web data acquisition for local businesses.
These findings provide the analytical dataset from which the research questions
posed in Chapter~1 will be systematically investigated in the following chapters.
